% =====================================================================
%  1 Datenmodell  (Aufgabenstellung: „Beschreibung der zentralen
%  Geschäftsobjekte der Anwendung und deren Beziehungen zueinander.
%  Textuell beschreiben und durch ein UML-Klassendiagramm visualisieren.“)
%  Fachliches Modell (Geschäftsobjekte), KEINE technischen Klassen.
%  Bewertung: „Sind die UML-Diagramme korrekt und fehlerfrei?“
% =====================================================================
\section{Datenmodell}
\TODO{Kurze Hinführung}

\subsection{Geschäftsobjekte}
\TODO{Je Geschäftsobjekt (z. B. Bereich, Eintrag, Ziel …): Zweck, wichtigste Attribute mit Datentyp und Pflicht/optional}

\subsection{Beziehungen und Klassendiagramm}
\TODO{Beziehungen und Multiplizitäten erläutern; UML-Klassendiagramm in docs/diagramme/02\_Klassendiagramm\_Datenmodell/ – Export als docs/diagramme/export/datenmodell.pdf}
\begin{figure}[H]
  \caption{Datenmodell als UML-Klassendiagramm}\label{abb:datenmodell}
  \centering
  \grafik{datenmodell}
  \quelle{Eigene Darstellung.}
\end{figure}
